\chapter{Requirements}

\section{Anforderungsprinzip}

Requirements beschreiben, was das System leisten oder einhalten muss. Technische Loesungen werden erst spaeter entschieden. Eine Aussage wie \enquote{Der Roboter verwendet einen bestimmten LiDAR} ist daher zunaechst keine Anforderung, sondern eine moegliche Architekturentscheidung. Diese Trennung entspricht der MBSE-Methodik des Lehrbuchs \cite[Kap.~6]{becker2026sysml}.

Jede Anforderung erhaelt:
\begin{itemize}
  \item eine stabile ID,
  \item einen kurzen Namen,
  \item einen Typ,
  \item eine Quelle oder ein Stakeholder Need,
  \item ein Akzeptanz- oder Nachweiskriterium,
  \item einen Status.
\end{itemize}

\section{Initiale Systemanforderungen}

\begin{longtable}{L{0.15\textwidth} L{0.16\textwidth} L{0.36\textwidth} L{0.18\textwidth}}
\toprule
ID & Typ & Requirement & Akzeptanzidee \\
\midrule
\modeltable{requirements_a1_initial}
\bottomrule
\end{longtable}

\section{MVP-relevante Anforderungen}

Fuer den ersten MVP werden zunaechst die Anforderungen priorisiert, die den technischen Kern aus Lokalisierung, Navigation, Umgebungserfassung, Hindernisbehandlung, Bewegungsausfuehrung und Statusmeldung pruefen.

\begin{longtable}{L{0.18\textwidth} L{0.16\textwidth} L{0.48\textwidth}}
\toprule
Requirement & MVP-Status & Begruendung \\
\midrule
\modeltable{mvp_a1}
\bottomrule
\end{longtable}

\section{Offene Quantifizierungen}

Folgende Anforderungen sind bewusst noch nicht numerisch festgelegt:
\begin{itemize}
  \item genaue Geschwindigkeitstoleranz und Fahrprofil fuer $0{,}5\,\mathrm{m/s}$,
  \item erlaubte Positionsabweichung,
  \item minimale Hindernisdetektionsdistanz,
  \item Bremsweg und Sicherheitsabstand,
  \item Nutzlast und Kistengeometrie,
  \item Batterieauslegung fuer das 2-h-Unterrichtsprofil,
  \item maximale Abmessungen des Roboters,
  \item Verfuegbarkeit und zulassige Fehlerrate.
\end{itemize}

Diese Werte werden nicht geraten, sondern in einer spaeteren Iteration aus Nutzungskontext, Safety-Konzept, mechanischem Aufbau und Testmoeglichkeiten abgeleitet.

\section{Ergaenzte Anforderungen aus Iteration A}

Aus den detaillierten Needs, Betriebsfaellen und dem Kartenkonzept ergeben sich folgende zusaetzliche oder praezisierte Anforderungen fuer die naechste Arbeitsfassung:

\begin{longtable}{L{0.15\textwidth} L{0.16\textwidth} L{0.36\textwidth} L{0.18\textwidth}}
\toprule
ID & Typ & Requirement & Akzeptanzidee \\
\midrule
\modeltable{requirements_a1_iteration_a}
\bottomrule
\end{longtable}

\section{Randbedingungen als Requirements-Quellen}

\begin{longtable}{L{0.15\textwidth} L{0.35\textwidth} L{0.32\textwidth}}
\toprule
Randbedingung & Bedeutung & Betroffene Requirements \\
\midrule
C-001 & Indoor-Betrieb & REQ-NAV-001, REQ-SAFE-001, REQ-MAP-001 \\
C-002 & Autonomer Fahrbetrieb & REQ-MIS-001, REQ-NAV-001, REQ-SAFE-001 \\
C-003 & Definierte Ziele & REQ-MIS-001, REQ-NAV-001 \\
C-004 & Vorhandene Grundkarte & REQ-MAP-001, REQ-LOC-001 \\
C-005 & Dynamische Umgebung & REQ-MAP-002, REQ-MAP-003, REQ-NAV-002, REQ-SAFE-001 \\
\bottomrule
\end{longtable}
